\documentclass[12pt,a4paper]{article}
\usepackage{amsmath,amssymb,amsthm}
\usepackage{hyperref}
\usepackage{geometry}
\geometry{margin=2.5cm}
\usepackage{enumitem}
\usepackage{booktabs}

\newtheorem{theorem}{Theorem}[section]
\newtheorem{lemma}[theorem]{Lemma}
\newtheorem{corollary}[theorem]{Corollary}
\newtheorem{proposition}[theorem]{Proposition}
\theoremstyle{definition}
\newtheorem{definition}[theorem]{Definition}
\theoremstyle{remark}
\newtheorem{remark}[theorem]{Remark}

\title{Black Hole Singularity Resolution in Recognition Science:\\
Discrete Ledger Extremal States}

\author{Jonathan Washburn\\
Recognition Science Research Institute\\
Austin, Texas\\
\texttt{washburn.jonathan@gmail.com}}

\date{March 2026}

\begin{document}
\maketitle

\begin{abstract}
We show that black hole singularities are dissolved in Recognition
Science (RS).  The discrete ledger has a minimum voxel size $\ell_0$
and minimum tick $\tau_0$, preventing infinite compression.  The
classical singularity ($r = 0$, $\rho = \infty$) is replaced by an
\emph{extremal ledger state}: a region where the recognition density
saturates at $\rho_{\max} = 1/\ell_0^3$ (one recognition event per
voxel).  The Planck density
$\rho_{\mathrm{Pl}} \approx 5.2 \times 10^{96}\;\mathrm{kg/m}^3$ is
the RS-predicted saturation scale.  Inside the horizon, the $J$-cost
curvature is bounded by $J_{\max} = J(\varphi^{n_{\max}})$, where
$n_{\max}$ is the maximum rung on the $\varphi$-ladder accessible at
the core.  The Penrose singularity theorem is not violated---its
mathematical conclusion (geodesic incompleteness) is reinterpreted:
geodesics terminate not at a singularity but at the extremal surface
where the discrete structure becomes manifest.
\end{abstract}

%% ============================================================
\section{Introduction}
%% ============================================================

The Penrose singularity theorem~\cite{Penrose1965} proves that, under
general assumptions (trapped surface, energy conditions, global
hyperbolicity), gravitational collapse produces geodesic
incompleteness---a ``singularity'' where curvature and density diverge.
This implies a breakdown of general relativity at the core of black
holes.

Proposed resolutions include loop quantum gravity (bounce
cosmology~\cite{Ashtekar2006}), string theory (fuzzball
conjecture~\cite{Mathur2005}), and noncommutative geometry.  Each
introduces new structures at the Planck scale.

Recognition Science (RS) provides a simpler resolution: the
discrete ledger has a built-in minimum length $\ell_0$ that
prevents infinite compression.  No new physics is
introduced---the discreteness that resolves the singularity is
the same discreteness that produces quantum mechanics (T2 in
the RS forcing chain).

RS derives all physics from the Recognition Composition Law
(RCL)~\cite{RS_Axioms}:
\begin{equation}
  J(xy) + J(x/y) = 2J(x)J(y) + 2J(x) + 2J(y),
\end{equation}
which uniquely determines the cost functional
$J(x) = \tfrac{1}{2}(x + x^{-1}) - 1$.  The forcing chain
(T0--T8) produces discreteness (T2: continuous
configurations cannot stabilize under $J$), the ledger
(T3: $J(x) = J(x^{-1})$ forces double-entry), and
$D = 3$ spatial dimensions (T8).

%% ============================================================
\section{The Discrete Ledger and Its Limits}
\label{sec:ledger}
%% ============================================================

\begin{definition}[Fundamental Scales]
The recognition ledger has:
\begin{itemize}
  \item Minimum length: $\ell_0 > 0$ (one voxel edge).
  \item Minimum time: $\tau_0 > 0$ (one tick).
  \item Minimum volume: $V_{\min} = \ell_0^3$.
\end{itemize}
\end{definition}

\begin{definition}[Maximum Recognition Density]
\label{def:rho-max}
The maximum density of recognition events is one event per voxel:
\begin{equation}
  \rho_{\max} = \frac{1}{\ell_0^3}.
\end{equation}
\end{definition}

\begin{theorem}[Density Bound]
\label{thm:density-bound}
For any physical configuration in the ledger:
$\rho \leq \rho_{\max}$.
\end{theorem}

\begin{proof}
Each voxel can host at most one recognition event per tick (the
Pauli-like exclusion on the ledger).  A voxel of volume $\ell_0^3$
contributes at most $1/\ell_0^3$ to the density.  Subdividing below
$\ell_0$ is impossible: the discrete lattice has no substructure.
\end{proof}

\begin{remark}[Status of the exclusion principle]
The ``Pauli-like exclusion'' (one recognition event per voxel per tick)
is a structural axiom of the ledger ontology, motivated by the double-entry
constraint $J(x) = J(x^{-1})$ (T3) which forces each ledger cell to
carry a unique, balanced recognition entry.  Multiple simultaneous
entries in a single voxel would require the same physical degree of
freedom to carry more than one value simultaneously---violating the
uniqueness of the recognition state.  A rigorous derivation of this
exclusion from the RCL alone is an identified open problem; the bound
$\rho \le \rho_{\max}$ is currently a structural postulate consistent
with the RS forcing chain.
\end{remark}

%% ============================================================
\section{The Classical Singularity}
\label{sec:classical}
%% ============================================================

In general relativity, the Schwarzschild solution has
\begin{equation}
  \rho(r) \to \infty \quad\text{as } r \to 0,
\end{equation}
and the Kretschmann scalar $K = R_{\mu\nu\rho\sigma}
R^{\mu\nu\rho\sigma} \sim r^{-6}$ diverges.  The geodesic
incompleteness proved by Penrose means that infalling worldlines
terminate in finite proper time.

\begin{remark}
The singularity theorem proves geodesic incompleteness under GR's
assumptions.  It does not prove that \emph{physical} density
diverges---that conclusion requires the additional assumption that
spacetime is a smooth continuum at all scales.  RS rejects this
assumption.
\end{remark}

%% ============================================================
\section{The Extremal Ledger State}
\label{sec:extremal}
%% ============================================================

\begin{definition}[Extremal State]
The extremal ledger state is the configuration in which every voxel
in the core region hosts a recognition event at the maximum
$\varphi$-rung $n_{\max}$:
\begin{equation}
  x_i = \varphi^{n_{\max}} \quad\text{for all voxels } i
  \text{ in the core}.
\end{equation}
The maximum rung $n_{\max}$ is set by the total mass-energy of the
black hole: $M_{\mathrm{BH}} = N_{\mathrm{core}} \cdot m_e \cdot
\varphi^{n_{\max}}$, where $N_{\mathrm{core}}$ is the number of
voxels in the core region.  For a solar-mass black hole,
$n_{\max} \approx 137$ (the fine-structure-constant rung); for a
stellar-mass BH, $n_{\max}$ is of this order.
\end{definition}

\begin{theorem}[Finite Core Cost]
\label{thm:finite-cost}
The $J$-cost at the core is bounded:
\begin{equation}
  J_{\max} = J(\varphi^{n_{\max}})
  = \frac{1}{2}\left(\varphi^{n_{\max}} + \varphi^{-n_{\max}}\right)
  - 1 < \infty.
\end{equation}
\end{theorem}

\begin{proof}
$\varphi^{n_{\max}}$ is a finite positive real number (the
$\varphi$-ladder has a maximum rung determined by the total mass-energy
of the black hole).  $J$ evaluated at any finite positive argument is
finite.
\end{proof}

\begin{theorem}[No Singularity]
\label{thm:no-singularity}
In Recognition Science, no physical quantity diverges at the center
of a black hole:
\begin{enumerate}[label=(\roman*)]
  \item Density: $\rho \leq \rho_{\max} = 1/\ell_0^3 < \infty$.
  \item Curvature: $K \leq K_{\max} \sim 1/\ell_0^4 < \infty$.
  \item $J$-cost: $J \leq J_{\max} < \infty$.
\end{enumerate}
\end{theorem}

\begin{proof}
All three bounds follow from the discreteness of the ledger.
(i)~Theorem~\ref{thm:density-bound}.
(ii)~Curvature is a second derivative of the metric, which on the
lattice has minimum step $\ell_0$; hence the maximum second difference
is $\sim 1/\ell_0^2$ per spatial direction, giving $K \sim 1/\ell_0^4$.
(iii)~Theorem~\ref{thm:finite-cost}.
\end{proof}

%% ============================================================
\section{Reinterpretation of the Penrose Theorem}
\label{sec:penrose}
%% ============================================================

\begin{theorem}[Geodesic Termination]
\label{thm:geodesic}
The Penrose singularity theorem's conclusion (geodesic incompleteness)
is reinterpreted in RS: infalling geodesics terminate not at a
singularity of infinite curvature but at the \emph{extremal surface}
where the continuous-geometry approximation breaks down and the
discrete ledger structure becomes manifest.
\end{theorem}

\begin{remark}
The extremal surface is located at
$r_{\mathrm{ext}} \approx \ell_0 \sim \ell_{\mathrm{Pl}}$
(the precise SI value of $\ell_0$ depends on the CODATA
calibration anchor $\hbar = E_{\mathrm{coh}} \cdot \tau_0$).
This is the RS analog of the ``Planck
star''~\cite{Rovelli2014} in loop quantum gravity, but
requires no new quantization procedure---the discreteness
is built into the ledger.
\end{remark}

%% ============================================================
\section{The Planck Density}
\label{sec:planck-density}
%% ============================================================

\begin{theorem}[Saturation Scale]
The maximum physical density in RS is
\begin{equation}
  \rho_{\max} = \frac{1}{\ell_0^3}
  \sim \rho_{\mathrm{Pl}},
\end{equation}
where $\rho_{\mathrm{Pl}} = c^5/(\hbar G^2) \approx
5.2 \times 10^{96}\;\mathrm{kg/m}^3$ is the Planck
density.  The precise numerical coefficient depends on
the calibration $\ell_0 \sim \ell_{\mathrm{Pl}}$, but
the saturation density is of order $10^{95}$--$10^{97}$
kg/m$^3$ regardless.
\end{theorem}

\begin{remark}
The saturation density being of order the Planck density
is consistent with the expectation that new physics (in
this case, the discrete ledger) becomes dominant near the
Planck scale.  This is analogous to the Hagedorn temperature
in string theory, where the density of states prevents
arbitrarily high temperatures.
\end{remark}

%% ============================================================
\section{Comparison with Other Approaches}
\label{sec:comparison}
%% ============================================================

\begin{center}
\renewcommand{\arraystretch}{1.3}
\begin{tabular}{@{}lcc@{}}
\toprule
\textbf{Approach} & \textbf{Resolution Mechanism} &
\textbf{Free Params} \\
\midrule
GR (classical) & Singularity (unresolved) & 0 \\
Loop QG & Quantum bounce & Barbero--Immirzi $\gamma$ \\
String theory & Fuzzball / stretched horizon & String scale \\
Noncommutative geom. & Smeared source & $\theta$ \\
\textbf{RS (ledger)} & \textbf{Extremal ledger state} &
\textbf{0} \\
\bottomrule
\end{tabular}
\end{center}

%% ============================================================
\section{Predictions}
\label{sec:predictions}
%% ============================================================

\begin{enumerate}
  \item \textbf{No trans-Planckian physics}: No experiment
  will detect length scales below $\ell_0 \sim \ell_P$.

  \item \textbf{Extremal remnant}: Black hole evaporation
  terminates at a Planck-mass remnant in the extremal
  state, not at zero mass.  The remnant has mass
  $M_{\mathrm{rem}} \sim M_{\mathrm{Pl}} \approx
  2.2 \times 10^{-8}\;\mathrm{kg}$ and entropy
  $S_{\mathrm{rem}} \sim O(1)$.

  \item \textbf{Gravitational wave echoes}: The extremal
  surface at $r_{\mathrm{ext}}$ may produce post-merger
  gravitational wave echoes.  The echo delay time is
  $\Delta t \sim R_s \ln(R_s/\ell_0)/c$, which for
  stellar-mass black holes gives $\Delta t \sim O(1)$~ms
  (in principle detectable by future gravitational wave
  observatories~\cite{Abedi2017}).

  \item \textbf{Finite core entropy}: The core has finite
  entropy $S_{\mathrm{core}} \leq k_R \cdot
  (r_{\mathrm{ext}}/\ell_0)^2$, bounded by the number
  of extremal voxels.
\end{enumerate}

%% ============================================================
\section{Conclusion}
\label{sec:conclusion}
%% ============================================================

Black hole singularities are dissolved by the discrete ledger.  The
classical singularity ($r = 0$, $\rho = \infty$) is replaced by an
extremal ledger state at $r \approx \ell_0$, where density, curvature,
and $J$-cost are all bounded and finite.  The Penrose theorem is not
violated but reinterpreted: geodesics terminate at the extremal surface
where the continuum approximation ceases to be valid.  The resolution
requires zero free parameters---it is a structural consequence of the
same discreteness that produces quantum mechanics.

%% ============================================================
\begin{thebibliography}{99}

\bibitem{RS_Axioms}
J.~Washburn, ``Recognition Science: Foundations,''
RS Axioms Document (2026).

\bibitem{Penrose1965}
R.~Penrose, ``Gravitational Collapse and Space-Time Singularities,''
Phys.\ Rev.\ Lett.\ \textbf{14}, 57 (1965).

\bibitem{Ashtekar2006}
A.~Ashtekar, T.~Pawlowski, and P.~Singh, ``Quantum Nature of the
Big Bang,'' Phys.\ Rev.\ Lett.\ \textbf{96}, 141301 (2006).

\bibitem{Mathur2005}
S.~D.~Mathur, ``The Fuzzball Proposal for Black Holes: An Elementary
Review,'' Fortschr.\ Phys.\ \textbf{53}, 793 (2005).

\bibitem{Rovelli2014}
C.~Rovelli and F.~Vidotto, ``Planck Stars,'' Int.\ J.\ Mod.\ Phys.\
D \textbf{23}, 1442026 (2014).

\bibitem{Abedi2017}
J.~Abedi, H.~Dykaar, and N.~Afshordi, ``Echoes from the Abyss:
Tentative Evidence for Planck-Scale Structure at Black Hole
Horizons,'' Phys.\ Rev.\ D \textbf{96}, 082004 (2017).

\bibitem{Hawking1976}
S.~W.~Hawking, ``Breakdown of Predictability in Gravitational
Collapse,'' Phys.\ Rev.\ D \textbf{14}, 2460 (1976).

\end{thebibliography}

\end{document}
